\chapter{Project Workplan and Budget}
\label{chap:planning}

\section{Chapter Overview}

This chapter records the zero direct-cash budget and the project schedule from
September 2026 to March 2027. The schedule is intentionally concise. Detailed
technical procedures and verification activities are described in
Chapters~\ref{chap:methodology} and \ref{chap:results}.

\section{Zero Direct-Cash Budget}
\label{sec:budget}

The project had a direct-cash budget of \textbf{KES 0}. No equipment,
commercial software, cloud compute, data, printing, travel, fabrication, or
testing service was purchased from project funds.

\begin{table}[H]
\centering
\caption{Direct-cash project budget}
\label{tab:cash-budget}
\small
\begin{tabularx}{\textwidth}{
    >{\raggedright\arraybackslash}p{0.23\textwidth}
    >{\raggedright\arraybackslash}X
    >{\centering\arraybackslash}p{0.15\textwidth}}
\toprule
\textbf{Category} & \textbf{Basis for zero direct expenditure} &
\textbf{Budget (KES)}\\
\midrule
Development hardware &
Existing personal or university computers were used. No computer or upgrade
was purchased for the project. & 0\\[0.35em]

CAD and simulation software &
Autodesk educational access was used. No commercial license was purchased
from project funds. & 0\\[0.35em]

Python and report software &
Python, project libraries, Git, and the LaTeX toolchain were available without
direct project charges. & 0\\[0.35em]

Cloud computing and storage &
All computation and working data remained local or used existing remote
access. No paid cloud service was purchased. & 0\\[0.35em]

Internet and electricity &
Existing personal or university services supported the work. No separate
project subscription or utility purchase was budgeted. & 0\\[0.35em]

Printing, travel, and communication &
The report was prepared for digital submission. No project-specific printing,
travel, or communication expense was charged. & 0\\[0.35em]

Prototype fabrication and testing &
No physical component or paid laboratory test was included in the completed
scope. & 0\\[0.35em]

Contingency &
The contingency consisted of schedule margin and scope control, not cash. & 0\\
\midrule
\textbf{Total direct-cash budget} & & \textbf{0}\\
\bottomrule
\end{tabularx}
\end{table}

The zero budget does not mean that the project used no resources. Existing
computers, Autodesk educational access, internet, electricity, supervision,
and student labor were in-kind inputs. Table~\ref{tab:in-kind} records them
without assigning invented monetary values.

\begin{table}[H]
\centering
\caption{In-kind resources used but not monetized}
\label{tab:in-kind}
\small
\begin{tabularx}{\textwidth}{
    >{\raggedright\arraybackslash}p{0.25\textwidth}
    >{\raggedright\arraybackslash}X
    >{\raggedright\arraybackslash}X}
\toprule
\textbf{Resource} & \textbf{Provision basis} & \textbf{Project use}\\
\midrule
Computers & Existing personal or university access &
Python development, tests, Fusion CAD, simulation, report compilation, and
demonstration.\\

Autodesk Fusion & Educational entitlement &
Native CAD generation, physical properties, STEP export, and manual Static
Stress preparation.\\

Open-source software & Public software licenses &
Optimization, testing, data records, dashboard, charts, version control, and
report production.\\

Internet and electricity & Existing personal or institutional access &
Research, documentation, source synchronization, and software operation.\\

Student labor and supervision & Academic project activity &
Engineering definition, implementation, review, testing, analysis,
documentation, and presentation preparation.\\
\bottomrule
\end{tabularx}
\end{table}

\section{Budget Control}
\label{sec:budget-control}

The architecture helped keep the direct cost at zero. Analytical equations
handled the large search, while manual FEA was reserved for a few candidates.
The project used the in-application Fusion API and local JSON files instead of
paid cloud automation. Open-source Python tools and digital reporting avoided
software and printing costs.

Prototype fabrication and paid testing were not treated as unnecessary. They
were placed outside the completed academic scope. A service-ready design would
need a separate budget for material certification, fabrication, inspection,
instrumentation, proof loading, and any required hydrodynamic testing.

\section{Workplan: September 2026 to March 2027}
\label{sec:workplan}

The project was organized over seven months. September and October established
the research basis, engineering definitions, and bracket proof of concept.
November and December developed the search, records, and multi-part
architecture. January and February focused on native CAD generation,
verification, validation preparation, and results. March was reserved for the
final report, demonstration, corrections, and submission.

The Fusion proof of concept was the main technical decision point. When the
public API did not provide dependable automated Static Stress control, the
project retained analytical optimization and native CAD, then changed the
solver step to a manual, checklist-based process. Figure~\ref{fig:gantt}
summarizes this sequence. Documentation and report writing overlap the
implementation work because assumptions and evidence were recorded throughout.

\begin{figure}[H]
\centering
\begin{ganttchart}[
    hgrid,
    vgrid,
    x unit=1.05cm,
    y unit title=0.62cm,
    y unit chart=0.55cm,
    title label font=\bfseries\scriptsize,
    bar label font=\scriptsize,
    group label font=\bfseries\scriptsize,
    milestone label font=\scriptsize,
    bar/.append style={fill=nemoblue!35,draw=nemoblue},
    group/.append style={fill=nemoteal!35,draw=nemoteal},
    milestone/.append style={fill=nemoorange,draw=nemoorange}
]{1}{7}
\gantttitle{2026}{4}\gantttitle{2027}{3}\\
\gantttitle{Sep}{1}
\gantttitle{Oct}{1}
\gantttitle{Nov}{1}
\gantttitle{Dec}{1}
\gantttitle{Jan}{1}
\gantttitle{Feb}{1}
\gantttitle{Mar}{1}\\

\ganttgroup{Research and definition}{1}{2}\\
\ganttbar{Scope and literature}{1}{2}\\
\ganttbar{Bracket model}{1}{2}\\

\ganttgroup{Implementation}{2}{5}\\
\ganttbar{Fusion proof of concept}{2}{3}\\
\ganttbar{Sampling and optimization}{3}{4}\\
\ganttbar{Padeye and stabilizer}{4}{5}\\
\ganttbar{Native CAD generators}{5}{5}\\

\ganttgroup{Verification and reporting}{5}{7}\\
\ganttbar{Tests and CAD sweeps}{5}{6}\\
\ganttbar{FEA preparation and results}{6}{6}\\
\ganttbar{Documentation and report}{1}{7}\\
\ganttbar{Final QA and demonstration}{7}{7}\\
\ganttmilestone{Submission}{7}
\end{ganttchart}
\caption{Portrait Gantt chart for September 2026 to March 2027}
\label{fig:gantt}
\end{figure}
